\begin{block}{Métodos}
    Para evaluar las capacidades de \boxi{beamerposter}, se utilizaron las siguientes métricas y metodologías:

    \begin{enumerate}
        \item \textbf{Tiempo de Elaboración}:
              \begin{itemize}
                  \item Se midió el tiempo total requerido para crear un póster desde cero utilizando \boxi{beamerposter}.
                  \item Se registraron los tiempos de compilación.
              \end{itemize}

        \item \textbf{Facilidad de Uso}:
              \begin{itemize}
                  \item Se evaluó la curva de aprendizaje para usuarios de diferentes niveles.
              \end{itemize}

        \item \textbf{Flexibilidad de Diseño}:
              \begin{itemize}
                  \item Se exploraron las opciones de personalización disponibles en \boxi{beamerposter}.
                  \item Se verificó la capacidad de implementar diseños complejos.
              \end{itemize}

        \item \textbf{Calidad del Resultado Final}:
              \begin{itemize}
                  \item Un panel de expertos evaluó la calidad visual del póster generado.
                  \item Se compararon aspectos como la resolución, alineación, consistencia tipográfica.
              \end{itemize}

        \item \textbf{Compatibilidad y Dependencias}:
              \begin{itemize}
                  \item Se revisó la compatibilidad de \boxi{beamerposter} con otros paquetes comunes de \LaTeX.
                  \item Se identificaron las dependencias adicionales requeridas.
              \end{itemize}

        \item \textbf{Documentación y Soporte}:
              \begin{itemize}
                  \item Se evaluó la calidad de la documentación oficial de \boxi{beamerposter}.
                  \item Se revisó la disponibilidad de tutoriales, ejemplos y comunidades de soporte en línea.
              \end{itemize}
    \end{enumerate}

    % Estas métricas permitieron realizar una evaluación sistemática y objetiva de \boxi{beamerposter}, enfocándose en aspectos clave como la eficiencia, la facilidad de uso y la calidad de los resultados.
\end{block}
